\section{Rule-Based Heuristic Baselines}
\label{sec:heuristic_baselines}
\noindent
In addition to sophisticated optimization-based (MPC) and learning-based (DRL) controllers, it is crucial to benchmark performance against simple, rule-based heuristic algorithms. These heuristics serve as fundamental baselines to quantify the value added by more intelligent strategies. The \texttt{heuristics.py} module implements several such baseline controllers.

\subsubsection{As-Fast-As-Possible}
This is the most straightforward and intuitive charging strategy. The As-Fast-As-Possible (AFAP) controller, implemented as \texttt{ChargeAsFastAsPossible}, commands all connected EVs to charge at their maximum permissible rate at all times, irrespective of electricity prices, grid constraints, or user needs beyond being ready by departure. Its objective is singular: minimize the time to reach full capacity. While simple, it often leads to undesirable outcomes such as high charging costs during peak price hours and contributing to grid congestion. A variation, \texttt{ChargeAsFastAsPossibleToDesiredCapacity}, stops charging once the user's specified desired SoC is met, rather than continuing to 100\%.

\subsubsection{As-Late-As-Possible}
The As-Late-As-Possible (ALAP) strategy, implemented in the \texttt{ChargeAsLateAsPossible} class, represents the opposite temporal extreme to AFAP. For each EV, it calculates the latest possible moment to initiate charging at maximum power to ensure the battery is full precisely at the scheduled departure time. This is achieved by determining the minimum number of time steps required to charge and deferring the start of the charging session until that point. This approach can shift load away from peak periods if those peaks occur early in the parking session, but it is inflexible and cannot respond to unexpected price drops or grid signals. It also carries the risk of failing to fully charge the vehicle if the departure time is earlier than expected. The \texttt{ChargeAsLateAsPossibleToDesiredCapacity} variant applies the same logic but targets the desired SoC.

\subsubsection{Round Robin Schedulers}
The \texttt{RoundRobin} family of heuristics attempts a more nuanced approach to meet a given power setpoint. The basic \texttt{RoundRobin} algorithm maintains a queue of connected EVs. To meet the power target, it calculates how many EVs need to be charged simultaneously (based on an average charging power) and activates that number of EVs from the front of the queue. In each time step, the EVs that were just charged are moved to the back of the queue, ensuring a fair distribution of charging opportunities over time.

\noindent
More advanced versions like \texttt{RoundRobin\_GF} (Grid-Friendly) are designed to track a power setpoint more precisely. Instead of activating a fixed number of EVs, it sequentially activates EVs and sums their potential power contribution until the setpoint is met or exceeded, providing a finer level of control.
